\section{动态规划证明、背包和区间 DP}

动态规划不是数组技巧，而是对子问题图的压缩求解。暴力搜索中，一个节点表示当前处理到哪里、已经做了哪些选择、剩余资源是什么。若不同路径到达同一个状态后，未来可选集合完全相同，就出现重叠子问题。DP 把这个状态命名并保存结果，再按依赖顺序计算。状态定义越稳定，代码越短；状态缺失一项，转移就会偷偷依赖历史细节。

线性 DP 常按最后一步思考。最大子数组和问以当前位置结尾的最佳后缀，打家劫舍问考虑前 i 个房屋的最大收益，股票题问每天结束时持有或不持有的最大收益。最后一步方法的价值在于把全局最优拆成互斥候选：当前元素单独开始还是接在前面，当前房屋偷还是不偷，今天买、卖还是休息。每个候选都必须对应一个已经计算好的状态。

二维字符串 DP 需要精确定义前缀。最长公共子序列的 dp[i][j] 表示两个前缀的答案，不是以某个字符结尾的答案。编辑距离的 dp[i][j] 表示把 word1 前 i 个字符变成 word2 前 j 个字符的最少操作。多出第零行和第零列不是为了方便而已，它们代表空串，给插入和删除操作提供统一边界。

背包题最容易背错循环顺序。0-1 背包中每个物品只能用一次，一维压缩时容量必须倒序，确保读取的是上一轮物品阶段的状态。完全背包中物品可无限使用，容量正序允许当前物品在同一轮被再次使用。组合计数外层遍历物品，排列计数外层遍历容量。循环顺序不是模板差异，而是状态语义差异。

目标和、分割等和子集和零钱兑换 II 是背包三连。分割等和是可达性，目标容量为总和一半；目标和通过代数变换把正负号选择转成选若干数凑出正集合和；零钱兑换 II 是完全背包组合计数。三题放在一起能训练四个问题：物品是否只能用一次，状态保存布尔值还是方案数，容量遍历正序还是倒序，答案是否受零元素影响。

区间 DP 的关键是按长度枚举。戳气球中，如果最后戳破区间内某个气球，那么它左右两侧已经分别被处理完，边界气球仍在；矩阵链乘法中，最后一次乘法把区间切成左右两段；回文相关区间题中，外层字符关系决定内部区间是否有效。区间 DP 的状态通常是 dp[left][right]，依赖更短区间，因此要先枚举区间长度。

空间压缩必须能解释旧值来自哪里。二维 LCS 压成一维时，dp[j] 在更新前是上一行同列，dp[j-1] 是当前行左侧，左上角需要额外变量保存。若解释不清这三个来源，空间压缩很容易写错。初学阶段先写二维表，确认状态和转移，再压缩空间。

实验建议：把 0-1 背包容量循环改成正序，用一个物品重量为一、价值为一、容量为二的用例观察它被重复使用。把零钱兑换 II 外层内层调换，观察组合数变成排列数。给 LCS 打印整个 DP 表，再手动从右下角回溯出一个公共子序列，理解表中每个值的来源。

\begin{principlebox}{本节复盘}
阅读本节后，请回到相邻章节的 LeetCode 案例，例如 LeetCode 875 爱吃香蕉、LeetCode 72 编辑距离、LeetCode 743 网络延迟时间、LeetCode 215 数组第 K 大，把暴力瓶颈、优化依据、状态不变量、C++ 容器成本和边界测试重新写一遍。能从这些具体题迁移到变体题，才算真正掌握。
\end{principlebox}
